\section{A Formal Model of Evidence Debt (RQ2--RQ4)}
\label{sec:model}

The definitions of \cref{sec:definitions} refer to reconstruction cost and decay
without structure. This section supplies the structure: operational history as a
graph, evidence as a fragmented projection of it, reconstruction as linkage
inference, and decay as survival processes over sources. The model is chosen so that
every parameter is either directly observable or estimable by the procedures of
\cref{sec:metrics,sec:field}, and so that the qualitative predictions tested in
\cref{sec:simulation} follow from it rather than from intuition.

\subsection{History, Evidence, and Fragmentation}
\label{sec:model-graph}

\textbf{Ground truth.} Operational history is a typed, acyclic \emph{event graph}
$\Hist = (V, L)$: nodes $V$ are events and artifacts of a change (intent record,
content revision, authorization act, build, artifact, deployment, runtime
configuration, outcome), and links $L \subseteq V \times V$ are the true causal and
derivational relations among them. A \emph{chain} is the subgraph of one logical
change. $\Hist$ is what actually happened; no one holds it directly.

\textbf{Evidence.} The corpus $\Corp(t)$ (\cref{def:corpus}) is a set of
\emph{records}, each held in a store $s \in S$, each testifying to a node and
carrying zero or more \emph{link keys}: identifiers intended to name adjacent nodes
(a ticket id in a commit message, a commit hash in a build record, a digest in a
deployment record)~\cite{sigelman2010}. The corpus thus determines an
\emph{observed graph}: nodes for surviving records, and an edge wherever a surviving
key names a surviving record. Fragmentation is the model's default: records of one
chain are distributed across stores with independent schemas, clocks, and retention
regimes (\cref{sec:amplifiers}).

\textbf{Deficit classes.} Relative to schema $\Sch$, a deficit
(\cref{def:deficit}) is one of: a \emph{missing node} (the record never existed or
was deleted), a \emph{missing key} (the record exists but the link key was never
propagated or was stripped), or a \emph{degraded attribute} (e.g., a null timestamp,
an unresolvable identity). The distinction is not pedantic: \cref{sec:simulation}
shows missing nodes and missing keys have different cost signatures---keys are
compensable by inference at effort-and-error cost; nodes, once every copy is gone,
are not compensable at all.

\subsection{Reconstruction as Linkage Inference}
\label{sec:model-linkage}

A query $q$ (\cref{def:rec}) requires recovering a subset of true links $L_q$. Where
keys survive, recovery is lookup. Where they do not, recovery is \emph{record
linkage}: deciding, from attributes (service, actor, temporal proximity, content
similarity), which surviving records refer to adjacent events---precisely the
inference problem formalized by Fellegi and Sunter~\cite{fellegi1969} and developed
in entity resolution~\cite{christen2012}. This identification imports the right
consequences. Linkage has \emph{effort}: candidate sets must be enumerated and
compared, and effort grows with the ambient density of similar records. Linkage has
\emph{error}: with imperfect discriminators, some accepted matches are false; false
links do not announce themselves, so the reconstructed history is silently wrong.
And linkage has \emph{limits}: when discriminating attributes are themselves
degraded (timestamps nulled, identities recycled), candidates become
indistinguishable and links unrecoverable even though both endpoint records
survive.

For a single required link $\ell$, let $\mathrm{paths}(\ell, t)$ be its surviving
\emph{recovery paths} at $t$: the direct key, redundant keys (a digest recorded in
two places), inferential linkage over attributes, and extra-corpus paths (asking the
engineer who did it). Each path $r$ has expected effort $c_r(t)$ and error rate
$\varepsilon_r(t)$. Under a fixed confidence-threshold policy, the effort-minimizing
acceptable path gives the link cost
\[
c_\ell(t) \;=\; \min \{\, c_r(t) : r \in \mathrm{paths}(\ell, t),\;
\varepsilon_r(t) \le \bar{\varepsilon} \,\},
\]
with $c_\ell(t)$ undefined---the link irrecoverable (\cref{def:irrec})---when the
set is empty. This path-level minimum is the cost of obtaining an acceptable answer,
not the always-finite attempt cost $\ERC$ of \cref{def:rec}. For an irrecoverable
link, $\ERC$ instead includes the finite search needed to exhaust the declared source
set and certify failure; $\mathsf I_q$ records the resulting physical loss state.
Policy abstention $\mathsf{Abst}_q$ can occur earlier and is not evidence that this
exhaustion boundary has been reached. For recoverable
links, query attempt cost aggregates the selected path costs (with shared-work
economies we do not model further). A returned link inconsistent with ground truth
sets $\Err_q$; the $\lambda_q$ term of \cref{def:debt} is therefore what makes
false reconstruction part of the debt quantity rather than merely a diagnostic
beside it. More generally, $B_q^{\star}$ in \cref{def:debt} minimizes the joint
effort, accepted-error, and abstention burden over admissible policies; it need
not choose the effort-minimizing path when a more expensive path reduces expected
consequence loss. $\ED^{\star}(t)$ compares optimal total burden with the ideal
corpus; $\ED_\rho(t)$ makes the same comparison for a fixed policy.

\subsection{Why Not $ED = f(M, F, R, U, C)$}
\label{sec:model-notformula}

A natural first proposal---and the one we expect readers to reach for---is a form
$ED = f(M, F, R, U, C)$ over missingness, fragmentation, retention risk,
uncertainty, and reconstruction cost. We considered and declined the form, keeping
the factors. The
reason is that the five quantities are not independent inputs to be aggregated; they
occupy different roles in the mechanism. Missingness ($M$) parameterizes the
deficit census. Fragmentation ($F$) and correlation uncertainty ($U$) shape the
\emph{cost and error of inferential paths}: fragmentation multiplies the stores to
be searched, and $U$---driven by ambient record density and discriminator
quality---sets false-match rates~\cite{fellegi1969}. Retention risk ($R$) is a
property of \emph{decay}: it determines which paths survive to time $t$
(\cref{sec:model-decay}). Reconstruction cost ($C$) is not an input at all; it is
the \emph{output} the others produce through the path structure. A scalar
aggregation of the five would assign weights to quantities whose interaction is
mediated by minimization over paths---and the minimization is what generates the
phenomena that matter (plateaus, cliffs, superadditivity; \cref{sec:model-props}).
Any composite index built for reporting purposes (\cref{met:edi}) must therefore be
understood as an estimator of \cref{def:debt}, not as its definition.

\subsection{Decay: Survival of Recovery Paths}
\label{sec:model-decay}

Each recovery path $r$ for a link has a \emph{survival function}
$\surv_r(\tau)$: the probability it remains usable $\tau$ after chain creation. The
empirically dominant forms are: \emph{step survival} for policy-driven retention
(a CI log is present until day $T_r$, then gone~\cite{nist80092,ghdocs,
cloudtrail}); \emph{event survival} for migrations and decommissionings (present
until an organizational event at an uncertain date); \emph{hazard survival} for
human sources (available until the holder departs, with turnover hazard
$\haz$~\cite{rigby2016}); and \emph{drift degradation} for inferential paths
(discriminators weaken as identities are recycled and schemas
evolve~\cite{parnas1994,lehman1980}). Interest (\cref{def:interest}) is the
consequence: as paths die, $c_\ell(t)$ jumps to the next-cheapest survivor, and
when the last dies, the link crosses into irrecoverability.

\subsection{Propositions}
\label{sec:model-props}

The following are analytic consequences of the path model; proofs are elementary
and stated to make the dependency structure explicit. Their empirical counterparts
are exercised in \cref{sec:simulation}, which distinguishes design-entailed
demonstrations from the falsifiable tests (notably the pairwise interaction test
of \cref{prop:superadd}).

\begin{proposition}[Effort plateau under a fixed acceptance policy]
\label{prop:plateau}
If a link retains at least one surviving path whose cost equals the original
cheapest cost, its expected reconstruction cost is unchanged by the loss of any
number of other paths. Consequently, estate-level cost as a function of a
degradation rate is flat while redundancy lasts and departs from flatness only as
cheapest paths begin to be destroyed. No plateau claim follows for total burden
unless a surviving path also preserves the error and consequence profile of the
previously burden-optimal policy.
\end{proposition}
\begin{proof}
Immediate from $c_\ell = \min$ over surviving paths: the minimum is invariant to
removal of non-minimal elements.
\end{proof}

\begin{proposition}[Superadditive interaction of deficit classes]
\label{prop:superadd}
Let two degradation classes each destroy a disjoint subset of paths such that
neither alone destroys any link's full path set. Then each alone adds bounded
effort cost and zero irrecoverability, while their combination can make the link
irrecoverable. The irrecoverability indicator is therefore strictly superadditive
for that link; no corresponding claim about the effort component or total $\ED$
follows without additional cost inequalities.
\end{proposition}
\begin{proof}
By construction, irrecoverability requires the surviving-path set to be empty
(\cref{def:irrec}). If class $A$ leaves $P \setminus A$ nonempty and class $B$
leaves $P \setminus B$ nonempty, but $A \cup B \supseteq P$, then
$\Irr_\ell(A)=\Irr_\ell(B)=0$ and $\Irr_\ell(A\cup B)=1$. Hence $1>0+0$.
For a sound exhaustive policy the combination also forces abstention. Under another
policy, physical irrecoverability may coexist with $W$ or $N$. In every case the
physical state is distinct from the policy decision and introduces no undefined or
infinite reconstruction cost.
\end{proof}

\begin{proposition}[Effort plateau and irrecoverability cliff]
\label{prop:steps}
Under the survival forms of \cref{sec:model-decay}, the expected finite
reconstruction-attempt cost of a fixed deficit is piecewise: constant
while the cheapest surviving path persists, jumping at retention boundaries and
organizational events, with smooth growth only where a hazard-governed human path
is the cheapest survivor. When the last answer-producing path dies, the remaining
cost is the finite effort of certifying irrecoverability. Under a sound exhaustive
policy, abstention then switches on and coincides with the physical loss state;
under a conservative policy it may switch on earlier, while an unsound or
aggressive policy may instead return $W$. The error and
abstention components have the same plateau-and-cliff form only when path error and
the acceptance rule remain constant between degradation events; in general they
change as discriminators degrade. Thus no constant $\rho$ makes total burden grow as
$e^{\rho \tau}$ in general.
\end{proposition}
\begin{proof}
Piecewise-constancy between path deaths follows from \cref{prop:plateau}; jumps at
step/event survival boundaries follow from minimization passing to the next
survivor, or to finite source-exhaustion effort when none remains; the only smooth
component is the expectation over hazard-distributed death times of a human path,
which yields growth bounded by the next attempt state's cost, not unbounded
compounding.
Error burden follows the same argument only under the stated between-event
constant-error assumption; otherwise its shape is determined by discriminator
degradation.
\end{proof}

\begin{proposition}[Bounded burden; mandatory irrecoverability boundary]
\label{prop:writeoff}
With finite penalty weights $\lambda_q$ and $\pi_q$, $\ED^{\star}(t)$ (and any
finite policy-specific $\ED_\rho(t)$) is bounded above by
$\sum_q w_q(\bar{C}_q + \max\{\lambda_q,\pi_q\})$, where $\bar{C}_q$ is the cost of $q$'s most
expensive reconstruction-attempt state. Once a link is irrecoverable, further decay
cannot restore an answer-producing path; the finite certification effort remains
bounded by $\bar{C}_q$. At physical irrecoverability, a sound exhaustive policy must
abstain; another policy may incur accepted-error burden instead. In all cases the
burden remains finite and does not grow without bound; it reaches a realized loss
cap. Financial debt has no
analogous physical write-off boundary imposed by the world rather than chosen by a
creditor.
\end{proposition}
\begin{proof}
Directly from \cref{def:debt}: reconstruction-attempt cost is finite by
\cref{def:rec} and bounded here by $\bar{C}_q$, while each probability difference
is at most one. Each query's contribution is therefore at most
$w_q(\bar{C}_q+\max\{\lambda_q,\pi_q\})$, because the default policy's accepted-
error and abstention outcomes are mutually exclusive.
\end{proof}

\label{sec:model-anti}
\Cref{prop:steps,prop:writeoff} are the model's verdict on the metaphor, obtained
before any experiment: evidence debt has principal and carrying cost, but its
effort and abstention exposure under sound exhaustive policies are stochastic,
discontinuous, and driven by third
parties' retention clocks rather than by any contracted rate; error exposure can
also drift continuously as discriminators weaken. Answerability terminates at a
physical write-off boundary rather than compounding indefinitely.
\Cref{sec:disc-metaphor} completes the
argument that the constructs stand without the metaphor.

\begin{proposition}[Measurement reduction]
\label{prop:reduction}
$\ED^{\star}(t)$ is determined by four input families: (i)~the deficit/corpus census
(observable, \cref{def:deficit}); (ii)~the reconstruction protocol family,
including admissible policies, confidence and finite-exhaustion rules, and
per-path effort/error calibrations; (iii)~the survival schedule of paths
(observable from retention policies, HR data, and migration calendars); and
(iv)~the declared workload and consequence weights $(w_q,\lambda_q,\pi_q)$.
\end{proposition}
\begin{proof}
By inspection of \cref{def:debt} through \cref{sec:model-linkage,sec:model-decay}:
the expectation is over path survival (iii), applied through the declared
reconstruction protocols (ii) to the corpus described by (i), and aggregated
under (iv).
\end{proof}

Two honesty notes bound what \cref{prop:reduction} claims. First, it is an
input-enumeration result, not a feasibility result: input (ii)---per-path cost and
error calibrations for future queries---is the empirically hard part, obtainable
only by reconstruction drills (\cref{sec:field}) and bounded, not supplied, by
linkage theory~\cite{fellegi1969}. Second, inputs (i) and (iv) are relative to
declared governance choices ($\Sch$; $w_q, \lambda_q, \pi_q$), so $\ED$ values are
comparable across estates only under shared reference declarations
(\cref{sec:metrics,sec:disc-relativity}). RQ2 is therefore answered as:
\emph{quantifiable relative to declared governance inputs, with the empirical
burden concentrated in cost calibration}---an affirmative with stated conditions,
not an unconditional one. \Cref{sec:metrics} turns the inputs into named metrics;
\cref{sec:sim-workload} instantiates the entire chain, computing $\ED_\rho(t)$ end to end in
a closed world where every input is known.
